\documentclass[12pt]{book}
\usepackage[utf8]{inputenc}
\usepackage[T1]{fontenc}
\usepackage{amsmath,amssymb,amsfonts}
\usepackage{amsthm}

\begin{document}

\setcounter{page}{349}
\section{5. Behavior with respect to certain base changes.}

In this section we show that the functor \(f^{\Delta}\) and the morphism \(Tr_{f}\) are compatible with certain base changes which take residual complexes into residual complexes.

Definition. A morphism \(f: X \to Y\) of locally noetherian
preschemes is residually if
a) \(f\) is flat
b) \(f\) is integral [EGA II 6.1.1], and
c) the fibres of \(f\) are Gorenstein preschemes

Examples.
i. An open immersion is residually stable.
2. A composition of residually stable morphisms is residually stable (use [V 9.6]).
3. If \(X\) and \(Y\) are the spectra of fields \(k \subseteq K\), then \(f\)
is residually stable if and only if \(K\) is algebraic over \(k\).
4. If \(f: X \to Y\) is residually stable, then every point \(x \in X\) is closed in its fibre, since \(f\) is integral. Thus the fibres are zero-dimensional Gorenstein preschemes.

\setcounter{page}{350}
\begin{proposition}
If \(f: X \to Y\) is a residually stable morphism, and if \(u: Y' \to Y\)
is a morphism of finite type, then \(g: X' = X \times_{Y} Y' \to Y'\)
is also residually stable.
\end{proposition}

\begin{proof}
Clearly \(g\) is flat and integral. To show that the fibres of \(g\) are Gorenstein, let \(y' \in Y'\),
let \(y=u(y')\), and consider the map of fibres \(v: X'_{y'} \to X_{y}\).
Now \(X_{y}\) is Gorenstein since \(f\) is residually stable.
The fibres of \(v\) are tensor products of fields, one of which is finitely generated (namely \(k(y') / k(y)\)), hence Gorenstein [V 9.5]. Hence \(X'_{y'}\) is Gorenstein [V 9.6].
\end{proof}

\begin{lemma}
Let \(A\) be a noetherian local ring, and let \(I\) be an \(A\)-module. Then \(I\) is an injective hull of the residue field \(k\) of \(A\) if and only
a) \(I\) has support at the closed point of \(\Spec A\),
b) \(\Hom_{A}(k, I) \cong k\),
c) There is a sequence of ideals \(\alpha_{1} \supseteq \alpha_{2} \supseteq \cdots\) of \(A\) which form a base for the \(\mathfrak{m}\)-adic topology, such that for all \(n\),
\[
\operatorname{length}\left(A / \alpha_{n}\right)
= \operatorname{length}\left(\Hom_{A}\left(A / \alpha_{n}, I\right)\right).
\]
\end{lemma}

\setcounter{page}{351}
\begin{proof}
If \(I\) is an injective hull of \(k\), then a),b),c) are immediate, since \(\Hom_{A}(\cdot, I)\) is a dualizing functor for
modules of finite length. In fact, c) holds for any ideal which is primary for \(\mathfrak{m}\).

Conversely, let \(I\) be an \(A\)-module satisfying a),b),c), and let \(J\) be an injective hull of \(k\).
By b), we can find an injection \(k \subseteq I\), and since \(J\) is injective, the natural map \(k \to J\) extends to a map \(\psi: I \to J\).
I claim \(\psi\) is injective. Indeed, given \(y \in I\), \(y \neq0\), we have \(\mathfrak{m}^{n} y=0\) for sufficiently large \(n\).
Let \(n\) be the least such integer. Then there is an \(x \in \mathfrak{m}^{n-1}\) such that \(xy \neq0\), but \(\mathfrak{m}(xy)=0\), so \(xy \in k\).
Hence \(\psi(xy) \neq0\), and \(\psi(xy)=x\cdot\psi(y)\), so \(\psi(y)\neq0\), and \(\psi\) is injective.

Now choose a sequence of ideals as in c). Then we note that
\[
\psi\left(\Hom\left(A / \alpha_{n}, I\right)\right) \subseteq \Hom\left(A / \alpha_{n}, J\right),
\]
and both have the same length, namely the length of \(A / \alpha_{n}\).
Therefore they are equal. But \(I\) and \(J\) are the union of the submodules of elements annihilated by \(\alpha_{n}\), so we see that \(\psi\) is also surjective.
Thus \(I\) is isomorphic to \(J\), and so is an injective hull of \(k\).
\end{proof}

\setcounter{page}{352}
\begin{proposition}
Let \(f: X \to Y\) be a residually stable morphism, and let \(K^{\cdot}\) be a residual complex on \(Y\). Then \(f^{*}(K^{\cdot})\) is a residual complex on \(X\).
\end{proposition}

\begin{proof}
Clearly \(f^{*}(K^{\cdot})\) is a complex of quasi-coherent sheaves with coherent cohomology. We have only to check that there is an isomorphism
\[
\sum_{p} f^{*}\left(K^{p}\right) \cong \sum_{x \in X} J(x) .
\]

Thus we reduce to the following local statement: For each \(x \in X\), let \(y=f(x)\),
let \(A=\sheaf{O}_{y}\) be the local ring of \(Y\),
let \(B=\sheaf{O}_{x}\) be the local ring of \(X\),
let \(I\) be an injective hull of \(k=k(y)\) over \(A\),
then \(I\otimes_{A} B\) is an injective hull of \(k'=k(x)\) over \(B\).

We apply the criteria of the lemma above.
Since \(I\) is an injective hull of \(k\) over \(A\), we have
a) \(I\) has support at the closed point of \(\Spec A\)
\[
b) \Hom_{A}(k, I) \cong k
\]
c) For each \(n\), where \(\mathfrak{m}\) is the maximal ideal of \(A\),
\[
\operatorname{length}\left(A / \mathfrak{m}^{n}\right)
= \operatorname{length}\left(\Hom_{A}\left(A / \mathfrak{m}^{n}, I\right)\right).
\]

\setcounter{page}{353}
Now since \(x\) is closed in its fibre, \(I\otimes_{A}\) has support at the closed point of \(\Spec B\).
Since \(B\) is flat over \(A\), we have
\[
\Hom_{B}\left(B / \mathfrak{m} B, I \otimes_{A} B\right) \cong B / \mathfrak{m} B .
\]
Therefore
\[
\Hom_{B}\left(k', I \otimes_{A} B\right) \cong \Hom_{B}\left(k', B / \mathfrak{m} B\right) \cong k',
\]
since \(B/\mathfrak{m}B\) is an Artinian Gorenstein ring.
Finally note for each \(n\),
\[
\operatorname{length}\left(B / \mathfrak{m}^{n} B\right)
= \operatorname{length}\left(\Hom_{B}\left(B / \mathfrak{m}^{n} B, I \otimes_{A} B\right)\right).
\]

Thus the conditions of the lemma are fulfilled, and \(I \otimes_{A} B\) is an injective hull of \(k'\) over \(B\).
\end{proposition}

\begin{proposition}
Let \(Y\) be the spectrum of a local Artin ring \(A\).
Then one can find a residually stable morphism \(u: Y' \to Y\)
where \(Y'\) is the spectrum of a local Artin ring \(A'\) with algebraically closed residue field.
\end{proposition}

\begin{proof}
By [EGA O777 10.3.1] one can find a local Artin ring \(A'\) and a local homomorphism \(A \to A'\) such that \(A'\) is flat over \(A\), \(\mathfrak{m}A'=\mathfrak{m}'\), and \(k' =\) the algebraic closure of \(k\).
Then \(Y'=\Spec A'\) will do.
Indeed, \(Y'\) is flat over \(Y\). It is integral since it is Artinian, and its residue field is algebraic over that of \(Y\).
The only fibre is a field, which is Gorenstein.
\end{proof}

\setcounter{page}{354}
Now we come to the behavior of \(f^{\Delta}\) and \(Tr_{f}\) under residually stable base change.
Let \(f: X \to Y\) be a morphism of finite type, and let \(u: Y' \to Y\) be a residually morphism.
Let \(X'=X \times_{Y} Y'\), and let \(u',f'\) be as shown.

If \(f\) is a finite morphism, there is a natural isomorphism
\[
f'^{Y} u^{*} \to u'^{*} f^{Y} \tag{V}
\]
derived from [III 6.3].

On the other hand if \(f\) smooth, an isomorphism
\[
f'^{Z} u^{*} \to u'^{*} f^{Z} \tag{VI}
\]
derived from [III 2.1].

\begin{theorem}
For every morphism \(f: X \to Y\) of finite type, and every residually stable morphism \(u: Y' \to Y\) (using the notations above) there is an isomorphism
\[
d_{u,f}: f'^{\Delta} u^{*} \to u'^{*} f^{\Delta}
\]
such that

\setcounter{page}{355}
1)If \(v: Y'' \to Y'\) is another residually stable morphism, then \(d_{uv,f}=d_{u,f}d_{v,f}\)

2) If \(f,g\) are two consecutive morphisms of finite type, then the isomorphisms \(d_{u,f}\) are compatible with \(c_{f,g}\)

3) If \(f\) is a finite morphism, then \(d_{u,f}\) compatible with \(\psi_{f}\) via (V)

4) If \(f\) is a smooth morphism, then \(d_{u,f}\) compatible with \(\varphi_{f}\) via (VI).
\end{theorem}

Suggestion of Proof. Show first that V and VI are compatible with the isomorphisms I,II,III and IV of \S2.
Define \(d_{u,f}\) for \(f\) finite or smooth using 3) and 4). Then define \(d_{u,f}\) for arbitrary \(f\) and check its properties, following the construction of \(f^{\Delta}\) given in \S3.

\begin{theorem}
Let \(u\) and \(f\) be as in the previous theorem. Then there is a commutative diagram
\end{theorem}

\setcounter{page}{356}
Proof. Show first that the analogous diagram of \(f^{Y}\) is valid for finite morphisms. Then follow the construction of \(Tr_{f}\) to show that it is true in general.

\end{document}